\section{Conclusion}
\label{submission:conclusion}

In this work, we introduced \textbf{Exclusive Parameter Merging (EPM)}, a training-free routing framework for merging sparse task-specific expert models under extreme domain conflict. EPM offers a robust, conceptually straightforward, and training-free alternative to highly parameterized test-time adaptation methods. Instead, EPM mitigates spatial weight-space interference through \textbf{Exclusive Parameter Allocation (EPA)}: primarily routing each weight coordinate to the task expert with the dominant scaled update, while retaining a background fraction of the standard blend to preserve coherence. To calibrate our task scales while avoiding the high-dimensional Overfitting-Optimizer Paradox, we developed \textbf{Task-Level Coefficient Tuning (TLC-Tune)}, which optimizes only $K$ global scaling factors using a zero-order (1+1) Evolution Strategy on a tiny offline validation split.

Our comprehensive evaluations on a Vision Transformer backbone across MNIST, FashionMNIST, CIFAR-10, and SVHN demonstrate that standard average-based merging methods collapse under severe weight-space conflict. In contrast, EPM successfully mitigates representational collapse, outperforming standard average-based merging methods (such as TIES-Merging and Task Arithmetic) by up to $\sim$25.6\% and $\sim$5.2\% absolute accuracy respectively under dense settings, and maintaining a significantly more balanced multi-task performance floor under moderate sparsity.

The broader takeaway of EPM is that direct coordinate-exclusive routing operators, when properly regularized, can resolve representational conflicts without relying on highly parameterized, layer-wise continuously tuned variables or auxiliary networks. This highlights that targeted weight-space design can be both structurally robust and highly performant. In future work, we plan to scale EPM to larger foundation models and multi-modal settings. Additionally, we aim to investigate the feasibility of adapting Soft-EPA to cross-architecture merging (where models do not share a pre-trained initialization or architectural structure) using optimal transport alignment techniques. We also plan to explore dynamic, layer-wise adaptation of the coherence retention factor $\gamma$, investigating whether lower values of $\gamma$ in shallow layers (to enforce coordinate exclusivity) combined with higher values of $\gamma$ in deeper layers can better preserve both local features and global activation manifold alignment, continuing to explore the potential of coordinate-level parameter composition in deep learning.
